\documentclass[11pt,a4paper]{article}

\usepackage[T1]{fontenc}
\usepackage[utf8]{inputenc}
\usepackage{lmodern}
\usepackage{amssymb,amsmath,amsthm,mathtools}
\usepackage[margin=2.5cm]{geometry}
\usepackage{parskip}
\usepackage[hidelinks]{hyperref}
\usepackage{microtype}

% Theorem environments (arXiv / journal style)
\theoremstyle{plain}
\newtheorem{theorem}{Theorem}[section]
\newtheorem{proposition}[theorem]{Proposition}
\newtheorem{lemma}[theorem]{Lemma}
\newtheorem{corollary}[theorem]{Corollary}

\theoremstyle{definition}
\newtheorem{definition}[theorem]{Definition}
\newtheorem{example}[theorem]{Example}
\newtheorem{remark}[theorem]{Remark}

% Notation
\newcommand{\Type}{\mathbf{Type}}
\newcommand{\STree}{\mathrm{STree}}
\newcommand{\SPos}{\mathrm{SPos}}

\title{Binary Proof Search Trees\\
  as a Directed Container}
\author{Extracted from a machine-checked development}
\date{}

\begin{document}
\maketitle

\begin{abstract}
Binary search trees form a directed container in which leaves are
terminal nodes (dead ends or successes) and internal nodes are branching
points. Shapes are tree skeletons; positions are paths from the root;
the directed structure is path concatenation. We verify the five
directed-container laws by structural recursion on paths. The induced
comonad extracts the root label and duplicates by replacing each label
with the full sub-search from that node, preserving dead ends at every
depth. A coKleisli morphism is a proof strategy that may inspect the
entire search tree. All stated results have been machine-checked.
\end{abstract}

\tableofcontents
\bigskip

%======================================================================
\section{Introduction}
%======================================================================

Proof search may be organised as a binary tree: at each branching point
one chooses among (here, two) alternatives; leaves record terminal
states---dead ends or successes. The distinction between dead end and
success lives in the labelling of nodes, not in the shape. Navigation of
such a tree is path concatenation: from a position one may restrict to
the subtree there, or concatenate a further path.

Directed containers make this precise. A directed container is a
container equipped with a root position, a subtree operation, and a
path-concatenation (shift) operation, satisfying five equational laws.
The associated polynomial functor becomes a comonad; coKleisli morphisms
are strategies that consume a whole labelled search tree.

\paragraph{What this note claims.}
In ordinary mathematical language: binary tree shapes with path
positions form a directed container. The five laws hold by structural
recursion on paths. Dead ends---leaves, whose only position is the
root---compose coherently at arbitrary depth: the search history is
self-consistent no matter where one zooms in.

%======================================================================
\section{Shapes and positions}
%======================================================================

\begin{definition}[Binary search tree shape]\label{def:stree}
A \textbf{binary search tree shape} is an element of the inductive type
$\STree$ generated by
\begin{itemize}
  \item $\mathsf{leaf}$ --- a terminal node;
  \item $\mathsf{node}(\ell, r)$ --- a branching point with left subtree
        $\ell$ and right subtree~$r$.
\end{itemize}
The distinction between dead end and success is not recorded in the
shape; it lives in any labelling of positions by a payload type~$X$.
\end{definition}

\begin{definition}[Position]\label{def:spos}
A \textbf{position} (path) in a shape $t \in \STree$ is an element of the
inductive family $\SPos(t)$ generated by
\begin{itemize}
  \item $\mathsf{here}$ --- the root of~$t$;
  \item $\mathsf{left}(p)$ --- descend into the left child of a node and
        continue along path~$p$;
  \item $\mathsf{right}(p)$ --- descend into the right child and continue
        along path~$p$.
\end{itemize}
For $t = \mathsf{leaf}$, the only position is $\mathsf{here}$. This is
what makes a leaf a dead end for further search moves.
\end{definition}

%======================================================================
\section{Subtree and path concatenation}
%======================================================================

\begin{definition}[Subtree]\label{def:sub}
The \textbf{subtree} operation
$\mathord{\downarrow} : (t : \STree) \to \SPos(t) \to \STree$
is defined by structural recursion on the path:
\begin{align}
  t \mathbin{\downarrow} \mathsf{here}
  &\;=\; t, \\
  \mathsf{node}(\ell, r) \mathbin{\downarrow} \mathsf{left}(p)
  &\;=\; \ell \mathbin{\downarrow} p, \\
  \mathsf{node}(\ell, r) \mathbin{\downarrow} \mathsf{right}(p)
  &\;=\; r \mathbin{\downarrow} p.
\end{align}
\end{definition}

\begin{definition}[Shift]\label{def:shift}
The \textbf{shift} (path concatenation)
\begin{equation}
  \oplus \;:\;
  (t : \STree) \to
  (p : \SPos(t)) \to
  \SPos\bigl(t \mathbin{\downarrow} p\bigr) \to
  \SPos(t)
\end{equation}
embeds a position in a subtree back into the original tree:
\begin{align}
  \mathsf{here} \oplus q
  &\;=\; q, \\
  \mathsf{left}(p) \oplus q
  &\;=\; \mathsf{left}(p \oplus q), \\
  \mathsf{right}(p) \oplus q
  &\;=\; \mathsf{right}(p \oplus q).
\end{align}
\end{definition}

%======================================================================
\section{The directed-container laws}
%======================================================================

A directed container consists of a container $(\mathrm{Shape}, \mathrm{Pos})$
together with operations
\begin{align}
  o &\;:\; \forall s.\; \mathrm{Pos}(s)
  &&\text{(root position)}, \\
  {\mathord{\downarrow}} &\;:\; \forall s.\; \mathrm{Pos}(s) \to \mathrm{Shape}
  &&\text{(subtree)}, \\
  {\oplus} &\;:\;
  \forall s.\;
  (p : \mathrm{Pos}(s)) \to
  \mathrm{Pos}(s \mathbin{\downarrow} p) \to
  \mathrm{Pos}(s)
  &&\text{(shift)},
\end{align}
satisfying five laws D1--D5. We take $\mathrm{Shape} = \STree$,
$\mathrm{Pos} = \SPos$, $o(t) = \mathsf{here}$, and the operations of
Definitions~\ref{def:sub} and~\ref{def:shift}.

\begin{proposition}[D1 and D2, definitional]\label{prop:d1-d2}
For every shape~$t$ and every position~$q$ in~$t$,
\begin{equation}
  t \mathbin{\downarrow} \mathsf{here} \;=\; t,
  \qquad
  \mathsf{here} \oplus q \;=\; q.
\end{equation}
\end{proposition}

\begin{proof}
By lean.
\end{proof}

\begin{proposition}[D3: right identity]\label{prop:d3}
For every shape~$t$ and every position~$p$ in~$t$,
\begin{equation}
  p \oplus \mathsf{here} \;=\; p.
\end{equation}
That is, concatenating the empty path on the right is the identity.
\end{proposition}

\begin{proof}
By lean.
\end{proof}

\begin{proposition}[D4: nesting of subtrees]\label{prop:d4}
For every shape~$t$, position~$p$ in~$t$, and position~$q$ in the subtree
at~$p$,
\begin{equation}
  t \mathbin{\downarrow} (p \oplus q)
  \;=\;
  \bigl(t \mathbin{\downarrow} p\bigr) \mathbin{\downarrow} q.
\end{equation}
\end{proposition}

\begin{proof}
By lean.
\end{proof}

\begin{proposition}[D5: associativity of shift]\label{prop:d5}
For every shape~$t$ and compatible positions $p$, $q$, $r$,
\begin{equation}
  (p \oplus q) \oplus r' \;=\; p \oplus (q \oplus r),
\end{equation}
where $r'$ is $r$ transported along the equality of fibres supplied
by~D4 (the position~$r$ lives over
$(t \mathbin{\downarrow} p) \mathbin{\downarrow} q$, while the left-hand
side requires it over $t \mathbin{\downarrow} (p \oplus q)$).
\end{proposition}

\begin{proof}
By lean.
\end{proof}

\begin{theorem}[Binary proof search as directed container]
\label{thm:proofsearch}
The data
\[
  \bigl(
    \STree,\;
    \SPos,\;
    o = \mathsf{here},\;
    {\mathord{\downarrow}},\;
    {\oplus}
  \bigr)
\]
form a directed container: laws D1--D5 hold.
\end{theorem}

\begin{proof}
By lean.
\end{proof}

%======================================================================
\section{The induced comonad}
%======================================================================

\begin{remark}[Extract and duplicate]
The directed container induces a comonad on the extension functor. For a
labelling $v : \SPos(t) \to X$,
\begin{align}
  \mathrm{extract}\,\langle t, v \rangle
  &\;=\; v(\mathsf{here}), \\[4pt]
  \mathrm{duplicate}\,\langle t, v \rangle
  &\;=\;
  \Bigl\langle
    t,\;
    p \mapsto
    \bigl\langle
      t \mathbin{\downarrow} p,\;
      q \mapsto v(p \oplus q)
    \bigr\rangle
  \Bigr\rangle.
\end{align}
Thus $\mathrm{extract}$ reads the label at the root, and
$\mathrm{duplicate}$ replaces the label at every node by the full
sub-search from that node---including all dead ends below it.
\end{remark}

\begin{remark}[Dead ends]
A dead end is a leaf: a position~$p$ at which
$\SPos(t \mathbin{\downarrow} p)$ is a singleton (only $\mathsf{here}$).
The comonad laws guarantee that dead ends compose coherently at
arbitrary depth: the search history is self-consistent no matter where
one zooms in.
\end{remark}

\begin{remark}[Proof strategies]
A coKleisli morphism $W(X) \to Y$ for the induced comonad~$W$ is a
\textbf{proof strategy}: a function that may inspect the entire search
tree (including dead ends) to decide what to produce.
\end{remark}

%======================================================================
\section{Discussion}
%======================================================================

\paragraph{What has been established.}
Binary search tree shapes with path positions form a directed container.
The five laws hold by structural recursion on paths; D1 and~D2 hold
definitionally as the base cases of that recursion. The induced comonad
navigates the search space while preserving dead ends at every depth.

\paragraph{Scope.}
The development treats strictly binary branching. The generalisation to
arbitrary finite arity (finitely branching rose trees) is a separate
construction with the same directed-container pattern.

\paragraph{Machine checking.}
Every definition and theorem statement above corresponds to a declaration
in a machine-checked formalisation. The phrase ``by lean'' indicates that
the proof has been verified by a proof assistant.

%======================================================================
\section*{Acknowledgements}
%======================================================================

The mathematical content follows the theory of directed containers of
Ahman, Chapman, and Uustalu. The presentation here is a prose rendering
of a formal development intended for human review of the theorem
statements.

\begin{thebibliography}{9}

\bibitem{ahman2012}
D.~Ahman, J.~Chapman, and T.~Uustalu.
\newblock When is a container a comonad?
\newblock In \emph{FoSSaCS}, volume~7213 of \emph{LNCS}, pages 74--88.
\newblock Springer, 2012.

\bibitem{abbott2005}
M.~Abbott, T.~Altenkirch, and N.~Ghani.
\newblock Containers: constructing strictly positive types.
\newblock \emph{Theoretical Computer Science}, 342(1):3--27, 2005.

\end{thebibliography}

\end{document}
